\section{Related Work}
\label{sec:related-work}

\textbf{Evaluation of LMs.}
Several recent works for evaluating LMs have either proposed a collection of mutually distinct tasks spanning across multiple domains~\citep{hendrycks2021measuring,liang2022holistic,srivastava2023imitation} or turned to the web as an interactive setting featuring tasks that require multiple steps to solve~\citep{yao2022webshop,zhou2023webarena,deng2023mind2web,liu2023agentbench}.
There are several drawbacks with such a ``potpourri" style setup.
First, each task tends to narrowly focus on one or a few skills, resulting in challenges that are typically too simple, pigeonhole the model into a reduced role, and do not provide models with the bandwidth to exercise their versatility or potentially demonstrate new abilities~\citep{srivastava2023imitation}.
Consequently, a model's performance on such task conglomerations may not yield actionable, deep insights regarding its capabilities and how to improve them~\citep{schlangen2019language,MartnezPlumed2021ResearchCD,bowman2021fix}.
\benchmark{} addresses these shortcomings, as our work demonstrates that it is significantly challenging, presents a wide range of possibilities for improving LMs to solve this task, and is easy to refresh over time with new task instances, each of which introduce novel, nuanced, and practical challenges.

\textbf{Code Generation Benchmarks.}
HumanEval~\citep{chen2021evaluating} is the current standard in a long-standing pursuit of synthesizing code from natural language descriptions~\citep{yu-etal-2018-spider,austin2021program,hendrycks2021measuring,li2022alphacode,zan2023large}.
In the past year, subsequent benchmarks have sought to augment HumanEval with extensions to different languages~\citep{cassano2022multiple,athiwaratkun2023multilingual,orlanski2023measuring}, variations in edit scope~\citep{yu2023codereval,du2023classeval}, similar but novel code completion tasks~\citep{muennighoff2023octopack}, and more testing~\citep{evalplus}.
Simultaneously, separate works have sought to introduce new coding paradigms~\citep{yin2022natural,yang2023intercode} or design library-specific problems~\citep{lai2022ds1000,zan2022cert}. 
Instead of partitioning problems into siloed datasets and curtailing them for simplicity's sake, \benchmark{}'s collection procedure transforms the source code with minimal post-processing, preserving a much broader set of challenges grounded in real-world software engineering beyond closed form completion, such as patch generation, reasoning over long contexts, navigating a codebase directory, and capturing dependency-based relationships across modules.

\textbf{ML for Software Engineering.}
To overcome traditional program analysis techniques that may not scale or incorporate natural language, one direction of current software engineering research is to use neural networks, including LMs, to automate real-world software development processes~\citep{google-didact,zheng2023towards,hou2023large}.
Use cases include automating commit generation~\citep{jung2021commitbert,liu2023commitbart}, PR review~\citep{Yang2016MSR,li2022automating,tufano2021automating}, bug localization~\citep{kim2019precise,chakraborty2018entropy}, testing~\citep{kang2023large,xia2023universal,wang2023software}, and program repair~\citep{gupta2017deepfix,allamanis2017learning,Monperrus_2018,jiang2018shaping,goues2019automated,gao2022program,dinh2023large,motwani2023better}.
Most relevant to \benchmark{} are works that have sought to apply LMs towards automated program repair~\citep{xia2022less,xia2023conversational,fan2023automated,sobania2023analysis}, guiding code editing with commits~\citep{chakraborty2021multimodal,zhang2022coditt5,fakhoury2023generating}.
However, none of the existing datasets~\citep{justJE2014defects,Karampatsis2019HowOD} present code context at the scale of \benchmark{}.
Moreover, \benchmark{} can be easily extended to new programming languages and repositories, and it provides a significantly more realistic and challenging arena to carry out experiments towards augmenting LMs with software engineering tools and practices.